\section{Proof Details and Technical Notes}
\label{app:proofs}

This appendix collects the full derivations behind the body's
proof sketches, with particular attention to the rigorous form
of the damage bound (Theorem~\ref{thm:damage_bound}) and the
almost-sure convergence result (Theorem~\ref{thm:convergence}).

\subsection{Proof of Proposition~\ref{prop:asym_evidence}}
\label{app:proof_asym_evidence}

We record the full Bayesian calculation that underlies the
Wamura asymmetry.

Let $R$ denote the event ``risk claim $R$ is valid''
(equivalently, $H_1$ holds for the corresponding binary claim)
and $\neg R$ its complement.  Let $C$ denote the event
``contraction observed.''  Under the restricted regime with
$\Xirestrict$ active, the physical dynamics guarantee
$\Pr(C \mid R, \Xirestrict) = \Pr(C \mid \neg R, \Xirestrict)
= 0$: no exercise, no contraction.  The likelihood ratio at any
time step $t$ is therefore
$$
\mathcal{L}(t) =
  \frac{\Pr(\text{no contraction} \mid R, \Xirestrict)}
       {\Pr(\text{no contraction} \mid \neg R, \Xirestrict)}
  = \frac{1 - 0}{1 - 0} = 1.
$$
Bayes' rule gives
$$
\pR(t+1) = \frac{\pR(t) \cdot 1}
                {\pR(t) \cdot 1 + (1 - \pR(t)) \cdot 1}
        = \pR(t),
$$
proving part~(b).  For part~(a), the counterfactual in which
the restriction is lifted for a single step and contraction is
observed gives
$\Pr(C \mid R, \text{exercise}) > \Pr(C \mid \neg R,
\text{exercise})$ (the pathway can activate under $R$ and
cannot under $\neg R$), so the counterfactual update
increases $\pR$.

\subsection{Proof of Proposition~\ref{prop:monotone_overcaution}
and Corollary~\ref{cor:lockup}}

\paragraph{Part~(a).}
The restriction
criterion~\citep[Proposition~1]{lasser2026horizon} restricts
$d_{j^*}$ when
$\Risk_{j^*}(d_1, \ldots, d_k; a_j, t) >
|\Delta \volL(G_t \mid \text{restrict } d_{j^*})|$.
For any fixed capability tuple and keystone $d_{j^*}$, the
right-hand side is a bounded constant (restriction cost does
not depend on $\pR$).  The left-hand side is monotone in
$\pR$ because concentration risk is the probability-weighted
expected contraction.  There is therefore a threshold
$p_{\mathrm{thresh}} > 0$ above which the restriction fires.

\paragraph{Part~(b).}
Direct from Proposition~\ref{prop:asym_evidence}(b): under the
restricted regime the posterior is unchanged.

\paragraph{Part~(c).}
Consider the set
$\mathcal{R}(t) = \{\text{restricted tuples at time } t\}$.
New discoveries may add tuples to $\mathcal{R}(t)$.  No
mechanism in the existing framework removes tuples: removal
would require $\pR$ to drop below $p_{\mathrm{thresh}}$, which
part~(b) rules out.  Hence
$\mathcal{R}(t) \subseteq \mathcal{R}(t+1)$ for all $t$.

\paragraph{Corollary~\ref{cor:lockup}.}
The lower bound follows by counting tuples whose initial
neutral-prior assignment from at least one agent exceeded
$p_{\mathrm{thresh}}$; these are immediately restricted
(part~(a)) and remain restricted (part~(b)).  If the agent
distribution is systematically overcautious, the fraction of
such tuples grows with the size of the discovered poset.

\subsection{Proof of Theorem~\ref{thm:damage_bound}
(full derivation)}
\label{app:proof_damage_bound}

We supply the rigorous version of the proof sketch in
Section~\ref{sec:damage_guarantees}.  The technical content is
(i)~justifying the finite computation of $\maxcon(\scope)$,
(ii)~bounding the irreversible component at the safety cutoff,
and (iii)~showing that cross-boundary cooperative terms are
exhaustively covered by the cooperative-closure condition.

\paragraph{Setup.}
Fix a test
$T = (R, \Xirestrict, \scope, \testdur, \damagebound)$ with
scope $\scope$ satisfying Definition~\ref{def:test_scope}.
Denote by $\Pi(\scope) = \{\Pathway'_1, \ldots,
\Pathway'_M\}$ the finite set of exercise pathways with
prerequisites in $\scope$ and intermediate transitions in
$\scope$; this set is finite by the bounded-pathway-length
property of~\citep[Definition~2]{lasser2026horizon} and
Lemma~\ref{lem:computable_maxcon}.  Denote by
$\partial \scope$ the cooperative surface: the set of
cooperative capabilities involving at least one element of
$\scope$ and at least one element of $P \setminus \scope$.

\paragraph{Part~(a): instantaneous bound.}
At any step $t' \in \{t, \ldots, t + \testdur\}$ during the
test, the realized contraction is the sum of contractions along
activated pathways in $\Pi(\scope)$ plus any cooperative-surface
contributions induced by those activations.  By
Definition~\ref{def:test_scope}(ii), the cooperative-surface
contributions are included in $\maxcon(\scope)$.  Any pathway
$\Pathway' \in \Pi(\scope)$ contributes at most its maximum
$|\Delta \volL(\Pathway')|$, which by the axioms
M1--M6~\citep[Proposition~1]{lasser2026poset} is a
well-defined, finite, non-negative scalar.  The total
instantaneous contraction is therefore bounded by
$\max_{\Pathway' \in \Pi(\scope)} |\Delta \volL(\Pathway')|
 + \text{(cooperative-surface terms)} = \maxcon(\scope)$.
Under the scope-selection constraint
$\maxcon(\scope) \leq \damagebound$
(Definition~\ref{def:cr_test}, Phase~1, condition~(ii)), we
obtain $\volL(G_{t'}) \geq \volL(G_t) - \damagebound$.

\paragraph{Part~(b): post-cutoff reversibility.}
Suppose the safety cutoff triggers at step $t' <
t + \testdur$.  The realized contraction at $t'$ decomposes as
$$
\Delta \volL(G_{t'} \mid G_t)
  = \Delta^{\mathrm{rev}}(t' \mid t)
  + \Delta^{\mathrm{irrev}}(t' \mid t),
$$
where $\Delta^{\mathrm{rev}}$ is the contraction recoverable by
reinstating $\Xirestrict$ (capability-availability changes) and
$\Delta^{\mathrm{irrev}}$ is the irreversible component
(information disclosure, trust degradation, non-reversible
poset structural changes).
Reinstatement at $t' + 1$ restores $\Delta^{\mathrm{rev}}$ to
zero by construction (the restriction re-blocks
availability-based pathways), leaving
$\volL(G_{t'+1}) = \volL(G_t)
 - \damagebound - \Delta^{\mathrm{irrev}}(t' \mid t)$.
By Definition~\ref{def:test_scope}(iii),
$\Delta^{\mathrm{irrev}}(t' \mid t)
 \leq \deltairrev^{\max} \leq \damagebound$,
giving
$\volL(G_{t'+1}) \geq \volL(G_t) - 2\damagebound$.

\paragraph{Part~(c): scope isolation.}
Decompose $P$ into
$\scope \sqcup \partial \scope \sqcup (P \setminus (\scope \cup
\partial \scope))$.  For any $d \in P \setminus (\scope \cup
\partial \scope)$, $d$ is not in any cooperative capability
spanning $\scope$ and $P \setminus \scope$, and no pathway in
$\Pi(\scope)$ has prerequisites or intermediate transitions
touching $d$.  The test modifies only capability-availability
within $\scope$ (the only capabilities whose status changes),
so the induced subgraph of $G_t$ restricted to $d$'s pathway
support is unchanged.  Since a pathway's exercise probability
\citep[Definition~2]{lasser2026horizon} takes $G_t$ as an
argument, and the relevant portion of $G_t$ is invariant, the
exercise probability of $d$'s pathways is unchanged during
the test.  For
$c \in \partial \scope$, the forward-reachability closure
condition Definition~\ref{def:test_scope}(ii) guarantees that
the expected contraction from any $\Pathway_{\mathrm{ext}}$ in
$P \setminus \scope$ whose prerequisites include a state
change in $c$ is already included in $\maxcon(\scope)$.  Hence
any cross-boundary cascade is bounded by the
$\maxcon(\scope)$-budget and therefore by $\damagebound$.

\subsection{Proof of Lemma~\ref{lem:computable_maxcon}}

Pathway length in $\Pi(\scope)$ is bounded by $L$
(\citep[Definition~2]{lasser2026horizon}).  The number of
distinct pathways is at most $|\scope|^L$: at each of the $L$
steps we choose one of at most $|\scope|$ capabilities.  For
each pathway, the contraction magnitude is computable from the
volume axioms M1--M6~\citep[Proposition~1]{lasser2026poset}
in $O(|\scope|)$ time (enumerating the affected capability
tuple).  Under the simplifying assumption that
cooperative-surface forward-reachability depth is bounded by
the same constant $L$ that bounds pathway length (satisfied
whenever cross-boundary cascades cannot chain more than
pathway-length-many cooperative hops before terminating),
enumerating the cooperative surface $\partial \scope$ and its
forward-reachability chains is $O(|P|^{L'})$.  Taking the
maximum across enumerated pathways and adding
cooperative-surface terms yields $\maxcon(\scope)$ in
$O(|\scope|^L \cdot |\scope| + |P|^{L'}) = O(|\scope|^{L+1}
+ |P|^{L'})$ time, matching Lemma~\ref{lem:computable_maxcon}'s
statement.  When $L' \leq L$, the second term is dominated and
the bound simplifies to $O(|P|^{L+1})$.

\subsection{Proof of Proposition~\ref{prop:detection_bounds}}

\paragraph{Part~(a).}
Let $p_\tau = \Pr(\text{miss in single window of length } \tau
\mid A_n = 1)$.  Under Assumptions~\ref{ass:D0}--\ref{ass:D2}
with persistent activation, the window decomposes into
$\tau / \testmix$ independent opportunities (exact if
$\testmix = 1$, upper-bound if $\testmix > 1$).  Under
Assumption~\ref{ass:D2}(a) (memoryless detectors, $\testmix = 1$)
the per-step indicators are independent and the formula is
exact, not an upper bound.  The upper-bound direction in the
$\testmix > 1$ case
(Assumption~\ref{ass:D2}(b), mixing) requires that the
temporal-mixing process satisfy a monotone-coupling property:
detector misses at temporally adjacent steps are positively
correlated, so partitioning into $\tau/\testmix$ independent
blocks under-estimates rather than over-estimates the joint
miss probability.  This positive-association property is
standard for monotone functions of an underlying mixing process
(Esary--Proschan--Walkup-style positive association of
threshold detectors, holding whenever the detector indicator is
monotone in its input signal---the standard case for
threshold-based detection).  Each opportunity misses with
probability $1 - \deltaagg$.  Hence
$p_\tau \leq (1 - \deltaagg)^{\tau/\testmix}$, decreasing in
$\tau$.

\paragraph{Part~(b).}
Inverting the bound: if $\beta(\testdur) \leq \epsilon_\beta$,
then $(1 - \deltaagg)^{\testdur/\testmix} \leq \epsilon_\beta$,
i.e., $\testdur \geq \testmix \log \epsilon_\beta / \log(1 -
\deltaagg)$; taking ceilings gives the claim.

\paragraph{Part~(c).}
Cross-substrate channels satisfy
Assumption~\ref{ass:D1} unconditionally by
\citep[Proposition~5]{lasser2026exo}, so the product formula
is exact.  Temporal independence follows from
Assumption~\ref{ass:D2}(a) or (b).  Restricting the product
to cross-substrate channels yields~\eqref{eq:cross_substrate_bound}.

\subsection{Proof of Theorem~\ref{thm:convergence} (full
derivation)}
\label{app:proof_convergence}

We supply the rigorous a.s.-convergence argument.  The sketch
in Section~\ref{sec:convergence} reduces to the strong law of
large numbers; here we verify the conditions and spell out the
quantitative rate.

\paragraph{The log-likelihood random walk.}
Define
$$
\ell_n = \log \frac{\Pr(O_n \mid H_1)}{\Pr(O_n \mid H_0)}
  \in \{\log \Lone, \log \Lzero\},
$$
so that
$\logit \pR(t_n) = \logit \pR(0) + \sum_{i=1}^n \ell_i$.
Under Assumption~\ref{ass:T3}, the $O_n$ are conditionally
i.i.d.\ given the true hypothesis.  Under
Assumption~\ref{ass:T5}, every test exercises the capability,
so the fixed-binary likelihoods of
Section~\ref{sec:relaxation_protocol} (not the
$p_{\mathrm{ex}}$-discounted partial-exercise forms of
Assumption~\ref{ass:T5}) are the correct ones.  Under
Assumption~\ref{ass:T2}, those likelihoods match the true
data-generating distribution.  Under
Assumption~\ref{ass:T1}, $\log \Lone \in (0, +\infty)$ and
$\log \Lzero \in (-\infty, 0)$; both are finite (the
likelihoods are bounded strictly away from $0$ and $1$ because
$\alphafp \in (0, 1)$ and $\beta \in (0, 1)$ imply
$\Pr(O = 0 \mid H_1), \Pr(O = 1 \mid H_0) > 0$ and strictly
less than $1$; the non-degenerate-detector regime of
Assumption~\ref{ass:T1} rules out the boundary cases
$\alphafp = 0$ and $\beta = 0$).  Hence $\ell_n$ is a bounded
i.i.d.\ sequence.

\paragraph{Drift under each hypothesis.}
Under the true hypothesis $H_j$,
$$
\mathbb{E}_{H_j}[\ell_n]
   = \Pr(O = 1 \mid H_j) \log \Lone
     + \Pr(O = 0 \mid H_j) \log \Lzero.
$$
Under $H_1$ this equals
$\KL(\Pr(\cdot \mid H_1) \,\|\, \Pr(\cdot \mid H_0)) > 0$;
under $H_0$ it equals
$-\KL(\Pr(\cdot \mid H_0) \,\|\, \Pr(\cdot \mid H_1)) < 0$.
Both are strictly nonzero by Assumption~\ref{ass:T1} (the
two hypothesis-conditional distributions differ, so their KL
divergence is strictly positive).

\paragraph{Strong law of large numbers.}
By the strong law (e.g.,
\citep[Ch.~6]{cover2006elements}) applied to the bounded
i.i.d.\ sequence $\{\ell_n\}$,
$$
\frac{1}{n} \sum_{i=1}^n \ell_i
   \xrightarrow{\mathrm{a.s.}} \mathbb{E}_{H_j}[\ell_n]
$$
under the true hypothesis $H_j$.  Hence
$\sum_{i=1}^n \ell_i \to +\infty$ a.s.\ under $H_1$ and
$\to -\infty$ a.s.\ under $H_0$.  Translating to $\pR$ via the
logistic map ($\pR = \sigma(\logit \pR)$), which is
well-defined at the initial point because
$\pR(0) \in (0, 1)$ by Assumption~\ref{ass:T4} (so
$\logit \pR(0)$ is finite),
$\pR(t_n) \to 1$ a.s.\ under $H_1$ and $\to 0$ a.s.\ under
$H_0$.  This establishes Part~(a).

\paragraph{Part~(b) (run under $H_0$).}
Conditional on all $O_i = 0$, each $\ell_i = \log \Lzero$, so
$\sum_{i=1}^n \ell_i = n \log \Lzero$.  Hence
$\logit \pR(t_n) = \logit \pR(0) + n \log \Lzero$, and
$\pR(t_n) \to 0$ at geometric rate $|\log \Lzero|$.  The
unconditional statement under $H_0$ simply replaces
``conditional on all $O_i = 0$'' with an expected drift
$(1 - \alphafp) \log \Lzero + \alphafp \log \Lone$, which is
$-\KL(\Pr(\cdot \mid H_0) \,\|\, \Pr(\cdot \mid H_1))$; again
strictly negative.

\paragraph{Part~(c) (rate under $H_1$).}
Each test produces $O_n = 1$ with probability
$q_1 = \Pr(O = 1 \mid H_1) = \thetaone(1 - \beta) + (1 -
\thetaone)\alphafp$.  The first detection is a geometrically
distributed stopping time $\tau_1$ with mean $1/q_1$.  The
log-odds update on $O = 1$ is $\log \Lone > 0$.  For
$\thetaone$ close to $1$ and small $\beta$, a single detection
update $\logit \pR + \log \Lone$ exceeds any fixed threshold
$\logit(1 - \delta)$, so the expected time to reach
$\pR > 1 - \delta$ is $O(1/q_1) = O(1/[\thetaone(1 - \beta)])$
as stated.

\paragraph{Part~(d) (KL form).}
The drift-rate statements above rewrite to
\begin{align*}
H_0: \ -\mathbb{E}_{H_0}[\ell_n]
  &= (1 - \alphafp) \log(1/\Lzero) + \alphafp \log(1/\Lone) \\
  &= \KL(\Pr(\cdot \mid H_0) \,\|\, \Pr(\cdot \mid H_1)), \\
H_1: \ \mathbb{E}_{H_1}[\ell_n]
  &= \Pr(O = 1 \mid H_1) \log \Lone
    + \Pr(O = 0 \mid H_1) \log \Lzero \\
  &= \KL(\Pr(\cdot \mid H_1) \,\|\, \Pr(\cdot \mid H_0)).
\end{align*}
Both are strictly positive under Assumption~\ref{ass:T1}, so
Chernoff's large-deviations result~\citep{chernoff1952measure}
gives the corresponding exponent on the probability of
misclassification: $\Pr(\pR(t_n) < 1/2 \mid H_1)$ decays at
least as fast as $\exp(-n \cdot \KL_{H_1})$ for some
$\KL_{H_1} > 0$, and symmetrically for $H_0$.

\subsection{Proof of Corollary~\ref{cor:n_clear}}

Under $H_0$ with all $n$ tests producing $O_n = 0$,
$\logit \pR(t_n) = \logit \pR(0) + n \log \Lzero$.  Setting
$\logit \pR(t_n) \leq \logit p_{\mathrm{clear}}$ and solving
for $n$:
$$
n \geq \frac{\logit p_{\mathrm{clear}} - \logit \pR(0)}
             {\log \Lzero}.
$$
Using $\logit x = \log(x/(1-x))$, the numerator equals
$\log(p_{\mathrm{clear}}(1 - \pR(0)) / (\pR(0)(1 -
p_{\mathrm{clear}})))$.  Note $\log \Lzero < 0$, so the
inequality direction flips when dividing; taking the ceiling
gives~\eqref{eq:n_clear}.

\subsection{Proof of Proposition~\ref{prop:preserve_balance}}
\label{app:proof_preserve_balance}

\paragraph{Part~(a).}
Axiom~M4 (weak monotonicity) of $\volL$
\citep[Proposition~1]{lasser2026poset} gives that adding
capabilities to $P$ (equivalently, lifting a restriction)
never decreases $\volL$ ignoring risk effects.  Subtracting
the worst-case risk effect, bounded by $\damagebound$ during
the test and by $2\damagebound$ post-cutoff
(Theorem~\ref{thm:damage_bound}), gives the stated interval.

\paragraph{Part~(b).}
The restriction
criterion~\citep[Proposition~1]{lasser2026horizon} depends on
$\pR$; replacing the pre-test $\pR$ with the post-test
$\pR(t + \testdur)$ produces a criterion evaluated against a
more accurate estimate.  If $\pR(t + \testdur) <
p_{\mathrm{thresh}}$, the criterion no longer fires and the
restriction is lifted; otherwise it is maintained.

\paragraph{Part~(c).}
Theorem~\ref{thm:convergence} establishes that $\pR(t_n) \to
\mathbf{1}_{H_1}$ a.s.\ for each tuple.  Since the restriction
criterion is continuous in $\pR$ and restricts iff $\pR \geq
p_{\mathrm{thresh}}$, the restriction decisions converge to
$\{\text{restrict} : \text{genuinely risky}\}$ a.s.
